\section{养老机构与长期照护中的性需求}

当老年人进入养老院、护理院或需要家人长期照护时，他们的性需求往往被彻底"抹去"——从公共话题中消失，甚至被默认为"不应该存在"。但性欲与对亲密、触摸、被爱的渴望并不会因为住进机构或需要照护而自动终止。这一节讨论一个很少被谈起的现实议题。

\subsection{被忽视的需求}

- \textbf{性欲持续存在}：研究显示，相当比例的老年人（包括 80 岁以上、部分失能者）仍保有性欲与性活动能力，性与亲密是生活质量、自我认同与尊严的重要组成。
- \textbf{需求转为"表达"而非消失}：即使因疾病、药物或体力无法进行传统性交，对拥抱、抚摸、依偎、言语亲密的需要依然强烈——"非性交的亲密"同样是性需求的一部分。
- \textbf{社会性沉默}：家属、护工、机构往往对老人的性与亲密避而不谈；一旦表达，常被当作"老不正经"或认知障碍的"异常行为"，造成二次伤害。

\subsection{机构与照护者的责任}

- \textbf{尊重自主与隐私}：在能力允许的前提下，尊重老年人对亲密关系与性表达的选择；提供可锁的私人空间与不被随意打扰的时段。
- \textbf{知情同意与权益保护}：对认知功能下降者，需要区分"自愿的亲密"与"被利用或被侵害"，建立保护机制，防止性侵与不当接触，而非简单地"一律禁止"。
- \textbf{培训与政策}：对护工进行性健康教育，明确"如何在不侵犯隐私的前提下保障安全"；机构应制定关于访客、伴侣留宿、亲密行为的管理规范。
- \textbf{疾病与用药评估}：性功能与欲望的变化可能与慢性病、药物（如抗抑郁药、降压药）有关，应由医生评估并可调整，而不是简单地"忍一忍"。

\subsection{给家属的建议}

- \textbf{把"性"重新理解为"亲密"与"尊严"}：与父母或长辈谈论时，重点放在"是否孤独、是否被尊重、是否需要陪伴"，而不是尴尬地回避。
- \textbf{支持而非审判}：若老人有新的伴侣或情感需求，理解其情感与身体需要，协助评估安全与健康，而不是斥责。
- \textbf{关注异常信号}：突然的、与人格不符的性言行，可能是认知障碍、脑病变、药物反应或谵妄的表现，需要就医评估。

\begin{tcolorbox}[colback=pink!5!white,colframe=red!50!black,title=核心原则]
住进养老机构、需要照护，都不等于"失去了作为人的性"。照护的目标是\textbf{在保障安全与尊严的前提下，尽可能保留自主与亲密}——既不回避，也不越界，而是把老年人当作有权表达情感与欲望的成年人。
\end{tcolorbox}
